\chapter{Dynamics and Railway State Interpretation}

\section{Motivation}

\begin{remark}
Geometric alignment alone is not sufficient to describe railway behaviour.
The motion of a vehicle depends on both curvature and cant, as well as
on operational parameters such as speed.
\end{remark}

\begin{remark}
This motivates a dynamic interpretation of alignment in which geometric
quantities are translated into physical effects.
\end{remark}

\section{Kinematic Basis}

\begin{definition}[Velocity]
Let $v>0$ denote the velocity of the vehicle along the alignment.
\end{definition}

\begin{definition}[Curvature-induced acceleration]
For a planar curve with curvature $\kappa(s)$, the lateral acceleration is
\[
a_\kappa(s) = v^2 \kappa(s).
\]
\end{definition}

\begin{remark}
This is the classical centripetal acceleration associated with curved motion.
\end{remark}

\section{Cant and Effective Acceleration}

\begin{definition}[Cant angle]
Let $\phi(s)$ denote the local roll angle induced by cant.
\end{definition}

\begin{definition}[Gravity component]
The gravitational acceleration projected onto the lateral direction is
\[
a_g(s) = g \sin\phi(s),
\]
where $g$ denotes gravitational acceleration.
\end{definition}

\begin{definition}[Effective lateral acceleration]
The effective lateral acceleration acting on the vehicle is
\[
a_{\mathrm{eff}}(s)
=
v^2 \kappa(s) - g \sin\phi(s).
\]
\end{definition}

\begin{remark}
This expression captures the fundamental coupling between curvature and
cant in railway engineering.
\end{remark}

\section{Equilibrium Condition}

\begin{definition}[Equilibrium cant]
A state of equilibrium is achieved when
\[
v^2 \kappa(s) = g \sin\phi(s).
\]
\end{definition}

\begin{remark}
In this case, the effective lateral acceleration vanishes:
\[
a_{\mathrm{eff}}(s) = 0.
\]
\end{remark}

\begin{remark}
This condition corresponds to the classical notion of balanced track,
where centrifugal and gravitational effects cancel.
\end{remark}

\section{Cant Deficiency and Excess}

\begin{definition}[Cant deficiency]
Cant deficiency is defined as
\[
\Delta(s)
=
v^2 \kappa(s) - g \sin\phi(s).
\]
\end{definition}

\begin{remark}
Positive values correspond to insufficient cant, negative values to
excess cant.
\end{remark}

\begin{remark}
Cant deficiency is a central design parameter in railway engineering,
as it directly influences passenger comfort and safety.
\end{remark}

\section{Dynamic Smoothness}

\begin{definition}[Jerk]
The rate of change of effective acceleration is given by
\[
j(s)
=
\frac{\dd}{\dd t} a_{\mathrm{eff}}(s).
\]
\end{definition}

\begin{remark}
Using $\dd t = \frac{1}{v} \dd s$, one obtains
\[
j(s)
=
v \frac{\dd}{\dd s} a_{\mathrm{eff}}(s).
\]
\end{remark}

\begin{proposition}
The jerk is given by
\[
j(s)
=
v^3 \kappa'(s) - g v \cos\phi(s)\,\phi'(s).
\]
\end{proposition}

\begin{remark}
This shows that both curvature variation and cant variation contribute
to dynamic effects.
\end{remark}

\section{Time Parameterization}

\begin{remark}
The alignment is parameterized by arc length \(s\), while dynamic
quantities evolve in time \(t\).
\end{remark}

\begin{definition}[Arc-length/time relation]
Assuming constant speed \(v>0\), we have
\[
\frac{\dd s}{\dd t} = v.
\]
\end{definition}

\begin{proposition}
For any function \(f(s)\), the time derivative is given by
\[
\frac{\dd f}{\dd t}
=
v \frac{\dd f}{\dd s}.
\]
\end{proposition}

\begin{remark}
This relation allows all spatial derivatives to be translated into
temporal dynamics.
\end{remark}

\section{Dynamic System Formulation}

\begin{definition}[State variables]
The railway state is described by
\[
x(s), y(s), \theta(s), \phi(s).
\]
\end{definition}

\begin{proposition}[Kinematic system]
The system satisfies
\[
\frac{\dd x}{\dd s} = \cos\theta(s),
\qquad
\frac{\dd y}{\dd s} = \sin\theta(s),
\]
\[
\frac{\dd \theta}{\dd s} = \kappa(s).
\]
\end{proposition}

\begin{remark}
These equations correspond to the geometric reconstruction of the curve.
\end{remark}

\begin{definition}[Extended dynamic state]
The extended state includes cant:
\[
(\gamma(s), \theta(s), \phi(s)).
\]
\end{definition}

\begin{remark}
The functions \(\kappa(s)\) and \(\phi(s)\) act as control variables.
\end{remark}

\section{Time-Based Dynamics}

\begin{proposition}
In time parameterization, the heading angle satisfies
\[
\frac{\dd \theta}{\dd t} = v \kappa(s).
\]
\end{proposition}

\begin{proposition}
The lateral acceleration is given by
\[
a_\kappa(t) = v^2 \kappa(s).
\]
\end{proposition}

\section{Higher-Order Dynamics}

\begin{definition}[Jerk]
The jerk is defined as
\[
j(t) = \frac{\dd}{\dd t} a_{\mathrm{eff}}(t).
\]
\end{definition}

\begin{proposition}
The jerk is given by
\[
j(t)
=
v^3 \kappa'(s)
-
g v \cos\phi(s)\,\phi'(s).
\]
\end{proposition}

\begin{proof}
Using
\[
a_{\mathrm{eff}} = v^2 \kappa(s) - g \sin\phi(s),
\]
we obtain
\[
\frac{\dd}{\dd t} a_{\mathrm{eff}}
=
v^2 \frac{\dd \kappa}{\dd t}
-
g \cos\phi(s) \frac{\dd \phi}{\dd t}.
\]
With
\[
\frac{\dd}{\dd t} = v \frac{\dd}{\dd s},
\]
this yields
\[
j(t)
=
v^3 \kappa'(s)
-
g v \cos\phi(s)\,\phi'(s).
\]
\end{proof}

\section{Control-Theoretic Interpretation}

\begin{definition}[Control system]
The railway alignment can be interpreted as a control system with
\[
\text{state: } (\gamma, \theta, \phi),
\qquad
\text{controls: } (\kappa, \phi').
\]
\end{definition}

\begin{remark}
The curvature \(\kappa(s)\) controls the geometric evolution,
while \(\phi(s)\) controls the dynamic balancing of forces.
\end{remark}

\begin{remark}
This formulation establishes a direct connection between alignment
design and optimal control theory.
\end{remark}

\section{Coupled Dynamics}

\begin{remark}
The effective acceleration depends on both curvature and cant:
\[
a_{\mathrm{eff}} = v^2 \kappa - g \sin\phi.
\]
\end{remark}

\begin{remark}
Thus, curvature and cant must be treated as a coupled system rather
than independent design variables.
\end{remark}

\begin{proposition}
Perfect equilibrium requires
\[
v^2 \kappa(s) = g \sin\phi(s).
\]
\end{proposition}

\section{Coupled State Interpretation}

\begin{remark}
The expressions above show that railway behaviour depends on the coupled
state variables
\[
(\kappa(s), \phi(s)).
\]
\end{remark}

\begin{remark}
A purely geometric interpretation based on $\kappa(s)$ alone is therefore
insufficient for describing the physical behaviour of the system.
\end{remark}

\section{Extension to Pose3}

\begin{remark}
Within the pose3 framework, the railway state is described by
\[
(x(s), y(s), z(s), \theta(s), \phi(s)).
\]
\end{remark}

\begin{remark}
This representation naturally integrates:
\begin{itemize}
\item geometry (position and direction),
\item vertical profile,
\item cant (through $\phi$),
\item and dynamic interpretation.
\end{itemize}
\end{remark}

\section{Relation to Optimization}

\begin{remark}
The dynamic quantities introduced above provide natural objective terms
for alignment optimization.
\end{remark}

\begin{remark}
Typical optimization criteria include:
\begin{itemize}
\item minimizing $\int a_{\mathrm{eff}}(s)^2 \,\dd s$,
\item minimizing $\int j(s)^2 \,\dd s$,
\item limiting peak values of cant deficiency.
\end{itemize}
\end{remark}

\section{Engineering Interpretation}

\begin{remark}
In practical railway design, acceptable ranges for cant, cant deficiency,
and jerk are prescribed by standards and operational guidelines.
\end{remark}

\begin{remark}
The mathematical formulation above provides a unified framework for
expressing these rules in analytic form.
\end{remark}

\section{Connection to Schwerpunkt-Trassierung}

\begin{remark}
The combined dependence on curvature and cant supports a
Schwerpunkt-oriented view of alignment.
\end{remark}

\begin{remark}
Rather than describing only a geometric path, alignment is interpreted as
a controlled dynamic state along the track.
\end{remark}

\section{Summary}

\begin{summary}
Railway alignment must be interpreted dynamically rather than purely
geometrically.

The coupled variables $\kappa(s)$ and $\phi(s)$ determine lateral
acceleration, comfort, and operational feasibility.

This provides the foundation for extending alignment modelling toward
state-based and optimization-based frameworks.
\end{summary}
